\newif\ifglobal

%%%%%%%%%%%%%%%%%%%%%%%
%%% Compile to view %%%
%%%%%%%%%%%%%%%%%%%%%%%

%%%%%%%%%%%%%%%%%%%%%%%%%%%%%%%%%%%%%%%%%%%%%%%%%%%
%%% Readme:                                     %%%
%%% https://www.overleaf.com/read/bwdjvfjksvjy/ %%%
%%%%%%%%%%%%%%%%%%%%%%%%%%%%%%%%%%%%%%%%%%%%%%%%%%%

% >>> M?: marks a module setting number or important notice

% >>> M4: Set {./relative-path} to external files for this module <<<
\def \modulefiles {./09-BOOTCTRL}
% To add an external file, use:
% (Replace '---' with actual filename)
% '\input{\modulefiles/tables/---.tex}' for tables
% '\includegraphics{\modulefiles/figures/---}' for figures
% '\subfile{\modulefiles/---}' for register files


% >>> M1: This module is in global project? <<<
% 'YES' - uncomment; 'NO' - comment out
\globaltrue

\ifglobal
    \documentclass[main\_\_CN.tex]{subfiles}
   
\else
    % Font "FandolSong-Regular" used by ctex does not have GSUB table.
    % It causes warnings that do not affect compilation and can be ignored.
    \PassOptionsToPackage{quiet}{fontspec}

    \documentclass[openany, 10pt]{book}
   
    % >>> M2: In ./00-shared/preamble-trm-module, also find
    %         settings for visibility of todo notes, watermark <<< 
    \usepackage{preamble-trm-module}


    % Enable Chinese typesetting
    \usepackage{ctex}

    % >>> M3: Temporary labels in Chinese
    %         you can add your own labels there <<<
    \myexternaldocument{./00-shared/temp-labels-cn}

\fi %\ifglobal


\setcounter{secnumdepth}{4}

% Define formatting of chapters in text
\titleformat{\chapter}
  [display]
  {\LARGE\bfseries}
  {\color{AllportsColor}第\, \thechapter\, 章}
  {1em}
  {\color{AllportsColor}}
\titlespacing*{\chapter}{0pt}{0pt}{20pt}

% Define headers
\usepackage{fancyhdr}
\usepackage{lipsum}
\renewcommand{\chaptermark}[1]{\markboth{第\, \thechapter\, 章 \  \ #1}{}}
\fancypagestyle{plain}{%
  \fancyhead{}%
  %\fancyhead[R]{\Navigation}
  \fancyhead[L]{%
    \nouppercase{%
      \ %
      \normalsize\leftmark
    }%
  }%
}
\pagestyle{plain}

% >>> M5: If this module needs any updates in
% './00-shared/preamble-trm-module.sty'
% (include additional package, etc.), contact your module PM
% or create the file './00-shared/changelog.tex' make the
% necessary updates yourself and document them in this file <<<

% Make text left-aligned and allow latex to break pages naturally, comment out for full justification
\raggedright
\raggedbottom


\begin{document}


% Sets language into which pre-defined bits of text are translated
\selectlanguage{chinese}

% Do not delete this block
\ifglobal
% Add the following front matter if
% this module is outside of global project
\else
    \InputIfFileExists{readme}{}{}
    \listoftodos
    {\let\clearpage\relax \listofdonetodos}
    \tableofcontents
    %\listoftables
    %\listoffigures
\fi


%%%%%%%%%%%%%%%%%%%%%%%%%
%%% TRM Module Begins %%%
%%%%%%%%%%%%%%%%%%%%%%%%%


\hypertarget{bootctrl}{}
\chapter{芯片 Boot 控制}
\label{mod:bootctrl}

\section{概述}

芯片启动过程以及芯片的某些功能是在上电或硬件复位时由 strapping 管脚和一些 eFuse 来确定的，具体包括以下功能：


\begin{itemize}
    \item \textbf{控制 Boot 模式}
    \begin{itemize}
        \item strapping 管脚：GPIO8 和 GPIO9
    \end{itemize}
    \item \textbf{SDIO 输入采样沿和输出驱动沿控制}
        \begin{itemize}
            \item Strapping 管脚：MTDI 和 MTMS
        \end{itemize}%
    \item \textbf{控制 ROM 日志输出到 UART0/USB}
    \begin{itemize}
        \item strapping 管脚: GPIO8
        \item eFuse 参数：EFUSE\_UART\_PRINT\_CONTROL 和 EFUSE\_DIS\_USB\_SERIAL\_JTAG\_ROM\_PRINT
    \end{itemize}
    \item \textbf{控制 JTAG 信号源}
    \begin{itemize}
        \item strapping 管脚：GPIO7
        \item eFuse 参数：
            EFUSE\_DIS\_PAD\_JTAG、EFUSE\_DIS\_USB\_JTAG 和 EFUSE\_JTAG\_SEL\_ENABLE
    \end{itemize}
\end{itemize}

上述 eFuse 位的默认值均为 0，也就是说没有烧写过。eFuse 只能烧写一次，一旦烧写为 1，便不能恢复为 0。有关烧写 eFuse 的信息，请参考章节 \ref{mod:efuse} \textit{\nameref{mod:efuse}}。

在上电复位和欠压复位过程中（请参考章节 \ref{mod:resclk} \textit{\nameref{mod:resclk}}），硬件将采样 strapping 管脚电平存储到锁存器中，并一直保持到芯片掉电或下一次芯片复位。strapping 管脚的锁存值可以通过软件从寄存器 \hyperref[fielddesc:GPIOSTRAPPING]{GPIO\_STRAPPING} 中读取。


\section{功能描述}
本小节主要介绍芯片复位时的功能以及控制该功能使用到的 strapping 管脚和 eFuse 组合模式。

\begin{importantlisting}\\
请使用本章节所介绍的组合，其它组合可能会导致不可控结果。
\end{importantlisting}

\subsection{默认配置}

GPIO9 默认连接内部上拉电阻。如果这一管脚没有外部连接或者连接的外部线路处于高阻抗状态，内部弱上拉将决定这一管脚输入电平的默认值，如表 \ref{tab:bootctrl-default-config-strapping-pins} 所示。

\begin{longtable}[c]{|p{1.5cm}|c|}
\caption{管脚默认上拉/下拉}
\label{tab:bootctrl-default-config-strapping-pins} \\
    \hline
    \rowcolor{lightgray}
  \textbf{管脚} & \textbf{默认值} \\ \hline
  \textbf{MTMS} & 浮空 \\ \hline
  \textbf{MTDI} & 浮空 \\ \hline
  \textbf{GPIO7} & 浮空 \\ \hline
  \textbf{GPIO8} & 浮空 \\ \hline
  \textbf{GPIO9} & 上拉 \\ \hline
\end{longtable}

如需改变 strapping 管脚的默认值，用户可以应用外部下拉/上拉电阻，或者应用主机 MCU 的 GPIO 来控制 \chipname{} 上电复位时的 strapping 管脚电平。复位释放后，strapping 管脚和普通管脚功能相同。


\subsection{Boot 模式控制} \label{sec:bootModeControl}


复位释放后，GPIO9 和 GPIO8 共同控制 Boot 模式。表 \ref{tab:bootctrl-boot-mode-control} 列出了 GPIO9 和 GPIO8 的 strapping 值及其对应的系统启动模式。


\begin{ThreePartTable}

\begin{TableNotes}
    \item[1] x：任何取值均不会对结果有影响，因此可忽略。
    \item[2] Joint Download Boot 模式：Joint Download Boot 模式下支持以下下载方式：
    \begin{itemize}
        \item USB-Serial-JTAG Download Boot
        \item UART Download Boot
        \item SDIO Slave 2.0 Download Boot
    \end{itemize}

\end{TableNotes}
\begin{longtable}[c]{|p{10cm}|c|c|}
\caption{系统启动模式}
\label{tab:bootctrl-boot-mode-control} \\
    \hline
    \rowcolor{lightgray}
  \textbf{启动模式} & \textbf{GPIO9} & \textbf{GPIO8} \\ \hline  
    \endfirsthead
    \insertTableNotes
    \endlastfoot
  SPI Boot 模式 & 1 & x\tnote{1}  \\ \hline 
  Joint Download Boot 模式\tnote{2} & 0 & 1   \\ \hline

\end{longtable}
\end{ThreePartTable}


在 \textbf{SPI Boot 模式}下，ROM 引导加载程序通过从 SPI flash 中读取程序来启动系统。SPI Boot 模式可进一步细分为以下两种启动方式：

\begin{itemize}
    \item \textbf{常规 flash 启动方式}：支持安全启动，ROM 引导加载程序将程序从 flash 加载到 SRAM 并执行。在大多数实际应用场景中，上述执行的程序多为二级引导程序，该二级引导程序将启动最终的应用程序。
    \item \textbf{直接启动方式}：不支持安全启动，程序直接从 flash 中运行。如需使能这一启动方式，请确保下载至 flash 的 bin 文件其前两个字为 0xaedb041d。详细的启动流程，见图 \ref{fig:boot_flow}。
\end{itemize}


在 \textbf{Joint Download Boot 模式}下，用户可通过 UART0 、USB 或 SDIO Slave 接口将二进制文件下载至 flash，或将二进制文件下载至 SRAM 并运行 SRAM 中的程序。



芯片启动的具体流程见图 \ref{fig:boot_flow}。


\begin{figure}[H]
    \centering
    {\includegraphics[width=1.2\textwidth]{\modulefiles/figures/boot_flow_250526.eps}}\\
    \caption{芯片启动流程}
    \label{fig:boot_flow}
\end{figure}

\begin{importantlisting}
    图中“Strapping Value”、“1xxx”和 “01xx”是 GPIO9、GPIO8 的组合值，见表 \ref{tab:bootctrl-boot-mode-control}。
\end{importantlisting}

下面几个 eFuse 可用于控制启动模式的具体行为：

\begin{itemize}
    \item \hyperref[fielddesc:EFUSEDISFORCEDOWNLOAD]{\textbf{EFUSE\_DIS\_FORCE\_DOWNLOAD}}
    \begin{itemize}
        \item 如果此 eFuse 设置为 0（默认），软件可通过设置 \hyperref[fielddesc:LPAONFORCEDOWNLOADBOOT]{LP\_AON\_FORCE\_DOWNLOAD\_BOOT}，触发 CPU 复位，将芯片启动模式强制从 SPI Boot 模式切换至 Joint Download Boot 模式。这种情况下，硬件会将寄存器 \hyperref[fielddesc:GPIOSTRAPPING]{GPIO\_STRAPPING}[3:2] 的值从 "1x" 覆盖为 "01"；
    
        \item 如果此 eFuse 设置为 1，则禁用 \hyperref[fielddesc:LPAONFORCEDOWNLOADBOOT]{LP\_AON\_FORCE\_DOWNLOAD\_BOOT}，寄存器 \hyperref[fielddesc:GPIOSTRAPPING]{GPIO\_STRAPPING} 的值不受影响 。 
    \end{itemize}
    
    \item \hyperref[fielddesc:EFUSEDISDOWNLOADMODE]{\textbf{EFUSE\_DIS\_DOWNLOAD\_MODE}}

    如果此 eFuse 设置为 1，则禁用 Joint Download Boot 模式，寄存器 \hyperref[fielddesc:GPIOSTRAPPING]{GPIO\_STRAPPING} 的值则不受 \hyperref[fielddesc:LPAONFORCEDOWNLOADBOOT]{LP\_AON\_FORCE\_DOWNLOAD\_BOOT} 的影响。
    
    \item \hyperref[fielddesc:EFUSEENABLESECURITYDOWNLOAD]{\textbf{EFUSE\_ENABLE\_SECURITY\_DOWNLOAD}}

    如果此 eFuse 设置为 1，则在 Joint Download Boot 模式下，只允许读取、写入和擦除明文 flash，不支持 SRAM 或寄存器操作。如已禁用 Joint Download Boot 模式，请忽略此 eFuse。
    \item \hyperref[fielddesc:EFUSEDISDIRECTBOOT]{\textbf{EFUSE\_DIS\_DIRECT\_BOOT}}

    如果此 eFuse 设置为 1，则禁用 Direct Boot 模式。
\end{itemize}

USB Serial/JTAG 控制器可将芯片从 SPI Boot 模式强制切换到 Joint Download Boot 模式，或从 Joint Download Boot 模式强制切换到 SPI Boot 模式。更多信息，请参考章节 \ref{mod:usbserialjtag} \textit{\nameref{mod:usbserialjtag}}。


\subsection{SDIO 输入采样沿和输出驱动沿控制}

MTMS 和 MTDI 管脚可用于调节 SDIO 输入采样沿和输出驱动沿。详见表 \ref{tab:bootctrl-sdio-input-output-edge-control} \textit{\nameref{tab:bootctrl-sdio-input-output-edge-control}}。

\begin{table}[H]
    \centering
    \caption{SDIO 输入采样沿/输出驱动沿控制}
    \label{tab:bootctrl-sdio-input-output-edge-control}

    \begin{threeparttable}
        \begin{tabular}{ | p{4.5cm} | c | c | }
        \hline
            \rowcolor{lightgray}
            \textbf{沿控制} & \textbf{MTMS\tnote{1}} & \textbf{MTDI\tnote{1}} \\\hline
            下降沿采样下降沿输出 & 0 & 0 \\\hline
            下降沿采样上升沿输出 & 0 & 1 \\\hline
            上升沿采样下降沿输出 & 1 & 0 \\\hline
            上升沿采样上升沿输出 & 1 & 1 \\\hline
         \end{tabular}

        \begin{tablenotes}
            \item[1] MTMS 和 MTDI 默认浮空，以上均非默认配置。
        \end{tablenotes}
     \end{threeparttable}
\end{table}


\subsection{ROM 代码日志打印控制}

系统在 SPI flash 启动模式下的 ROM 引导阶段，GPIO8、\hyperref[regdesc:LPAONSTORENREG]{LP\_AON\_STORE4\_REG}[0] 与 \hyperref[fielddesc:EFUSEUARTPRINTCONTROL]{EFUSE\_UART\_PRINT\_CONTROL} 一起控制 ROM 代码日志打印。



\begin{table}[H]
    \centering
    \caption{ROM 代码日志打印控制}
    \label{tab:bootctrl-rom-code-printing-control}
    \begin{threeparttable}
    \begin{tabular}{ |p{2.8cm}<{\centering}|p{3cm}<{\centering}|p{1.3cm}<{\centering}|p{5cm}|}
    \hline
        \rowcolor{lightgray}
    \textbf{\hyperref[regdesc:LPAONSTORENREG]{LP\_AON\newline\_STORE4\_REG}[0]} &   \textbf{\hyperref[fielddesc:EFUSEUARTPRINTCONTROL]{EFUSE\_UART\newline\_PRINT\_CONTROL}} &   \textbf{GPIO8}      & \textbf{ROM 代码日志打印} \\\hline
    \multirow{6}*{0}   & 0(0b00) &  x\tnote{1}  & ROM 代码日志始终打印至 UART0\\\cline{2-4}
     & \multirow{2}*{1(0b01)}   & 0     & 使能 \\\cline{3-4}
      & & 1       & 关闭 \\\cline{2-4}
     & \multirow{2}*{2(0b10)}   & 0     & 关闭 \\\cline{3-4}
      & & 1       & 使能 \\\cline{2-4}
      & 3(0b11) &  x  & 关闭\\\hline
     1 & x & x & 关闭\\\hline
        \end{tabular}

        \begin{tablenotes}
            \item[1] x: 任何取值均不会对结果有影响，因此可忽略。
        \end{tablenotes}
     \end{threeparttable}
    \end{table}

ROM 代码日志上电默认打印至 UART0 和 USB Serial/JTAG 控制器，可通过置位\hyperref[fielddesc:EFUSEDISUSBSERIALJTAGROMPRINT]{EFUSE\_DIS\_USB\_SERIAL\_JTAG\_ROM\_PRINT} 禁用 USB Serial/JTAG 控制器打印。

注意：在 \hyperref[fielddesc:EFUSEDISUSBSERIALJTAGROMPRINT]{EFUSE\_DIS\_USB\_SERIAL\_JTAG\_ROM\_PRINT} 设置为 0，即可打印至 USB 的情况下，如果 USB Serial/JTAG 控制器已被禁用，则 ROM 代码将无法打印到 USB Serial/JTAG 控制器。


\subsection{JTAG 信号源控制}

在系统启动早期阶段，GPIO7 可用于控制 JTAG 信号源。该管脚没有内部上下拉电阻，strapping 的值必须由不处于高阻抗状态的外部电路控制。

GPIO7 与 \hyperref[fielddesc:EFUSEDISPADJTAG]{EFUSE\_DIS\_PAD\_JTAG}、\hyperref[fielddesc:EFUSEDISUSBJTAG]{EFUSE\_DIS\_USB\_JTAG} 和 \hyperref[fielddesc:EFUSEJTAGSELENABLE]{EFUSE\_JTAG\_SEL\_ENABLE} 共同控制 JTAG 信号源，见表 \ref{tab:bootctrl-jtag-signal-src-control}。


\begin{table}[H]
    \centering
    \caption{JTAG 信号源控制}
    \label{tab:bootctrl-jtag-signal-src-control}
    \begin{threeparttable}
    \begin{tabular}{ | C{2cm} | C{2cm} | C{2.2cm} | C{1.5cm} | l | }
    \hline

    \rowcolor{lightgray}
    \textbf{\hyperref[fielddesc:EFUSEDISPADJTAG]{EFUSE\_DIS\newline\_PAD\_JTAG}} & \textbf{\hyperref[fielddesc:EFUSEDISUSBJTAG]{EFUSE\_DIS\newline\_USB\_JTAG}} & \textbf{\hyperref[fielddesc:EFUSEJTAGSELENABLE]{EFUSE\_JTAG\newline\_SEL\_ENABLE}} & \textbf{GPIO7} & \textbf{JTAG 信号源} \\\hline

    \multirow{3}*{0} & \multirow{3}*{0} & 0 &  x\tnote{1}  & \multirow{2}*{USB 串口/JTAG 控制器} \\\cline{3-4}
      &   & \multirow{2}*{1}   & 1     & \\\cline{4-5}
      &   &                    & 0     & \multirow{3}*{JTAG 管脚 MTDI、MTCK、MTMS 和 MTDO} \\\cline{1-4} 
    0 & x & x                  & x     & \\\cline{1-4}
    0 & 1 & x                  & x     &                                                 \\\hline
    1 & 0 & x                  & x     & USB 串口/JTAG 控制器 \\\hline
    1 & x & x                  & x     & \multirow{2}*{JTAG 关闭} \\\cline{1-4}
    1 & 1 & x                  & x     &                         \\\hline
        \end{tabular}

        \begin{tablenotes}
            \item[1] x：x 表示该值被忽略，任何取值不影响该状态。
        \end{tablenotes}
    \end{threeparttable}
\end{table}

\end{document}
